\chapter{Einleitung}

Ziel des Projekts war der Aufbau eines durchgängigen Demonstrators für industrielle
Bildverarbeitung und Machine Learning. Der erste Projektteil behandelt die klassische
Bildverarbeitung: ein Basler-Stereo-Kamerasystem wird kalibriert, Bildpaare werden
rektifiziert, daraus wird eine Disparitätskarte berechnet und anschließend eine farbige
3D-Punktwolke exportiert. Der zweite Projektteil ergänzt diese geometrische Verarbeitung um
Machine Learning: ein eigenes YOLO-Modell erkennt Schreibgeräte in den Rohbildern der
Kameras.

Die beiden Teile sind bewusst getrennt aufgebaut. Die Bildverarbeitung erzeugt metrische
3D-Informationen und hängt stark von Kamerageometrie, Beleuchtung, Textur und
Disparitätsparametern ab. Das Machine-Learning-Modul liefert semantische Information in Form
von Klassen, Konfidenzen und Bounding Boxes. Zusammen bilden sie die Grundlage für eine
spätere semantische 3D-Auswertung, bei der erkannte Objekte direkt mit Punkten im Raum
verknüpft werden könnten.

\section{Aufgabenstellung}

Aus zwei Kamerabildern soll eine verwertbare Punktwolke entstehen. Dafür werden zunächst die
intrinsischen Parameter beider Kameras sowie die relative Lage der Kameras zueinander
bestimmt. Danach werden neue Stereo-Szenen aufgenommen oder aus dem Projektordner geladen,
rektifiziert und mit einem Stereo-Matching-Verfahren verarbeitet. Die Tiefeninformation
entsteht aus der Disparität, also der horizontalen Verschiebung korrespondierender Bildpunkte.

Parallel dazu soll ein Detektor für drei Stiftklassen trainiert werden. Die Klassen sind
\textit{Fineliner}, \textit{Marker} und \textit{Pencil}. Der Detektor arbeitet auf den
unveränderten Rohbildern und ist damit unabhängig von der Rektifizierung und der
Punktwolkenberechnung.

\section{Projektziele}

Im Projekt wurden folgende Ziele umgesetzt:

\begin{itemize}
    \item Mono- und Stereo-Kalibrierung mit einem ChArUco-Board,
    \item Speicherung wiederverwendbarer Kalibrierparameter,
    \item Aufnahme und Wiederverwendung von Stereo-Bildpaaren,
    \item Rektifizierung, Disparitätsschätzung und 3D-Reprojektion,
    \item Export roher, geschnittener und bereinigter Punktwolken im PLY-Format,
    \item Integration einer YOLO-Inferenz auf linken und rechten Rohbildern,
    \item Erstellung und Training eines eigenen YOLO-Datensatzes für Schreibgeräte,
    \item Dokumentation der erzeugten Artefakte, Grenzen und nächsten Schritte.
\end{itemize}

\section{Abgrenzung}

Der aktuelle Stand ist ein funktionsfähiger Workflow und keine vollständig validierte
Messlösung. Die Punktwolkenqualität wurde visuell und über Plausibilitätswerte geprüft; eine
systematische Messunsicherheitsanalyse mit Referenzkörpern wurde noch nicht durchgeführt.
Auch die Objekterkennung ist bisher als 2D-Detektor integriert. Eine robuste Fusion von
Bounding Boxes und 3D-Punkten ist als nächster Entwicklungsschritt vorgesehen.
